\section{Metodologia}

\subsection{Conjunto de dados}

O conjunto de dados utilizado neste estudo é composto por 624 tarefas de programação dos conjuntos CEJava, CEPython e HumanEval. Para cada tarefa, foram geradas cinco implementações por diferentes modelos de linguagem, resultando em 3.120 implementações de código.

Cada implementação possui uma classificação manual das alucinações presentes no código, de acordo com a taxonomia proposta por Liu et al. A partir dessas classificações, foram identificados 12 tipos distintos de alucinação.

Além disso, cada implementação foi analisada individualmente pelo SonarQube, sendo registradas as regras violadas pelo código. Foram identificadas 71 regras distintas considerando o identificador completo da regra, incluindo a linguagem à qual pertence, como, por exemplo, \texttt{java:S1854} e \texttt{python:S1854}.

Regras Java e Python com o mesmo identificador numérico foram mantidas separadas, uma vez que pertencem a analisadores específicos de cada linguagem e não se assumiu equivalência entre elas.

Considerando os 12 tipos de alucinação e as 71 regras do SonarQube, foram definidos inicialmente

\[
12 \times 71 = 852
\]

pares alucinação--regra para investigação.

\subsection{Preparação dos dados}

\subsubsection{Binarização das variáveis}

Para permitir a análise da relação entre os tipos de alucinação e as regras do SonarQube, ambas as variáveis foram representadas por indicadores binários.

Para cada tipo de alucinação $X$, definiu-se:

\[
X =
\begin{cases}
1, & \text{se a implementação apresenta a alucinação;}\\
0, & \text{caso contrário.}
\end{cases}
\]

Analogamente, para cada regra do SonarQube $Y$:

\[
Y =
\begin{cases}
1, & \text{se a implementação apresenta uma violação da regra;}\\
0, & \text{caso contrário.}
\end{cases}
\]

Dessa forma, foi construída uma matriz no nível da implementação contendo uma linha para cada código gerado, além de indicadores binários para os 12 tipos de alucinação e para as 71 regras do SonarQube.

\subsubsection{Elegibilidade das implementações por linguagem}

Como as regras analisadas pelo SonarQube são específicas da linguagem, cada regra foi avaliada somente no subconjunto de implementações correspondente à sua linguagem.

Para regras identificadas pelo prefixo \texttt{java:}, foram consideradas as implementações Java provenientes do conjunto CEJava. Para regras identificadas pelo prefixo \texttt{python:}, foram consideradas as implementações Python provenientes dos conjuntos CEPython e HumanEval.

Esse subconjunto é denominado, neste estudo, \textit{conjunto de implementações elegíveis} para a regra analisada.

Essa restrição é necessária porque uma implementação escrita em Python, por exemplo, não deve ser interpretada como uma observação com ausência de uma violação Java. Nesse caso, a regra simplesmente não pertence ao conjunto de regras aplicáveis àquela implementação.

\subsubsection{Níveis de análise}

A análise foi estruturada em dois níveis: implementação e tarefa.

No nível da implementação, cada um dos 3.120 códigos gerados ocupa uma linha da matriz de dados e é analisado separadamente. A presença da alucinação $X$ e da regra $Y$ refere-se, portanto, à mesma implementação.

No nível da tarefa, as cinco implementações correspondentes à mesma
tarefa são agregadas em uma única unidade de análise. Para cada tipo
de alucinação e regra do SonarQube, uma tarefa recebe valor 1 quando
pelo menos uma de suas cinco implementações apresenta a característica
analisada, e valor 0 quando nenhuma implementação a apresenta.

Consequentemente, a presença simultânea de uma alucinação $X$ e de uma
regra $Y$ no nível da tarefa não implica que ambas tenham ocorrido na
mesma implementação. Ela indica que, entre as implementações associadas
àquela tarefa, pelo menos uma apresentou $X$ e pelo menos uma apresentou
$Y$.

Assim, no nível da tarefa, a ocorrência de uma alucinação e de uma regra não implica necessariamente que ambas tenham ocorrido na mesma implementação.

\subsection{Procedimento de análise exploratória}

Para caracterizar a frequência marginal das variáveis, considerou-se
$n_X$ o número de implementações que apresentam
a alucinação $X$, e $n_Y$ o número de
implementações que apresentam a violação $Y$.
A Tabela~\ref{tab:distribuicao_frequencias_marginais} apresenta a
distribuição dessas frequências entre os pares avaliáveis. 

\begin{table}[ht]
    \centering
    \caption{Distribuição das frequências marginais dos pares 
    alucinação--regra avaliáveis.}
    \label{tab:distribuicao_frequencias_marginais}
    
    \begin{tabular}{lrr}
        \hline
        \textbf{Estatística} 
        & $\boldsymbol{n_X}$ 
        & $\boldsymbol{n_Y}$ \\
        \hline
        
        Média              & 59,54  & 13,46 \\
        1º quartil (25\%)  & 3      & 1     \\
        Mediana (50\%)     & 26     & 3     \\
        3º quartil (75\%)  & 129    & 9     \\
        90º percentil      & 166,6  & 26    \\
        95º percentil      & 172    & 61    \\
        Máximo             & 291    & 371   \\
        
        \hline
    \end{tabular}
\end{table}\label{tab1_esparcidade}

Observa-se que as violações do SonarQube apresentam frequências
particularmente baixas. A mediana de $n_Y$ foi igual a 3, indicando
que, em metade dos pares avaliáveis, a respectiva regra do SonarQube
foi observada em no máximo três implementações elegíveis. Além disso,
o primeiro quartil de $n_Y$ foi igual a 1, mostrando a elevada
concentração de regras com baixa frequência de ocorrência.

Como pode ser observado na Tabela~\ref{tab:distribuicao_frequencias_marginais}, há elevada esparsidade observada nas combinações entre tipos de alucinação e regras do SonarQube e, devido a esta característica, a análise concentra-se na caracterização dos padrões observados nos dados, sem utilizar testes de significância como critério para determinar a existência de associações.

A análise considera conjuntamente a ocorrência das variáveis, a direção e a magnitude das diferenças observadas e a quantidade de observações que sustentam cada padrão.

\subsubsection{Construção das tabelas de contingência}

Para cada par formado por um tipo de alucinação $X$ e uma regra do
SonarQube $Y$, foi construída uma tabela de contingência $2\times2$,
conforme apresentado na Tabela~\ref{tab:contingencia}.

\begin{table}[ht]
    \centering
    \caption{Estrutura da tabela de contingência para cada par
    alucinação--regra.}
    \label{tab:contingencia}
    \begin{tabular}{lcc}
        \hline
        & $Y=1$ & $Y=0$ \\ 
        \hline
        $X=1$ 
        & $a$: $X$ e $Y$ presentes
        & $b$: apenas $X$ presente \\

        $X=0$
        & $c$: apenas $Y$ presente
        & $d$: $X$ e $Y$ ausentes \\
        \hline
    \end{tabular}
\end{table}

A partir das células da tabela de contingência, foram calculadas as
frequências marginais:

\[
n_{X=1}=a+b,
\qquad
n_{X=0}=c+d,
\]

\[
n_{Y=1}=a+c,
\qquad
n_{Y=0}=b+d.
\]

Essas frequências são utilizadas tanto para caracterizar a ocorrência
das variáveis quanto para contextualizar as medidas exploratórias
calculadas posteriormente.

\subsubsection{Caracterização da esparsidade}

Antes da interpretação das medidas de associação, foi avaliada a
esparsidade das tabelas de contingência.

Foram examinadas as frequências marginais de $X$ e $Y$, o número de
ocorrências conjuntas e a presença de células com frequência zero.

Não foi adotado um número mínimo de coocorrências como critério
automático de exclusão. A ausência de coocorrência não implica
necessariamente ausência de um padrão de associação, pois $X$ e $Y$
podem ocorrer individualmente com frequência e, ainda assim, não
ocorrer conjuntamente.

Os pares foram, portanto, classificados estruturalmente nas seguintes
categorias:

\begin{description}

    \item[$X$ ausente:]
    não existem implementações com $X$ no conjunto elegível para a
    regra $Y$, impossibilitando a comparação entre implementações com
    e sem $X$;

    \item[Sem coocorrência:]
    $X$ e $Y$ ocorrem no conjunto elegível, mas nenhuma implementação
    apresenta ambos simultaneamente;

    \item[Coocorrência completa:]
    existe pelo menos uma ocorrência conjunta de $X$ e $Y$, e todas
    as quatro células da tabela de contingência possuem frequência
    diferente de zero;

    \item[Coocorrência com célula zero:]
    existe pelo menos uma ocorrência conjunta de $X$ e $Y$, mas uma
    das células comparativas da tabela apresenta frequência zero.

\end{description}

Essa classificação permite distinguir a ausência de informação para
comparação da ausência efetivamente observada de coocorrência.

\subsubsection{Diferença de proporções}

Para cada par avaliável, foi calculada a proporção de implementações com a violação $Y$ entre aquelas que apresentaram a alucinação $X$:

\[
P(Y=1\mid X=1)
=
\frac{a}{a+b},
\]

e a proporção correspondente entre implementações sem a alucinação:

\[
P(Y=1\mid X=0)
=
\frac{c}{c+d}.
\]

A diferença entre essas proporções foi definida como:

\[
\Delta =
P(Y=1\mid X=1)
-
P(Y=1\mid X=0).
\]

Valores positivos de $\Delta$ indicam que, no conjunto observado, a regra $Y$ é proporcionalmente mais frequente entre implementações com a alucinação $X$. Valores negativos indicam que $Y$ é proporcionalmente mais frequente entre implementações sem $X$. Valores próximos de zero indicam proporções semelhantes entre os dois grupos.

A diferença de proporções é utilizada como medida descritiva da direção e da magnitude do padrão observado, e não como evidência inferencial de associação na população.

\subsubsection{Odds ratio}

Como medida complementar de associação, foi considerado o \textit{odds ratio} (OR):

\[
OR=\frac{ad}{bc}.
\]

Quando todas as células possuem frequência diferente de zero, valores de $OR>1$ indicam maior \textit{odds} de ocorrência de $Y$ entre implementações com $X$, enquanto valores de $OR<1$ indicam menor \textit{odds}. Um valor de $OR=1$ corresponde à igualdade das \textit{odds} entre os grupos.

Entretanto, devido à ocorrência frequente de células com valor zero, o OR bruto não é utilizado isoladamente para ordenar ou classificar os pares. Valores iguais a zero ou que tendem ao infinito são interpretados conjuntamente com a configuração da tabela de contingência e com as frequências que lhes deram origem.

Em particular, uma medida de associação extrema baseada em poucas observações não é interpretada da mesma maneira que uma medida de magnitude semelhante sustentada por um número maior de ocorrências.

\subsubsection{Suporte observacional}

A magnitude das diferenças observadas é analisada conjuntamente com o suporte observacional de cada par.

Para isso, são consideradas principalmente as frequências:

\[
n_X=n_{X=1}=a+b
\]

e

\[
n_Y=n_{Y=1}=a+c.
\]

Essa avaliação é necessária porque diferenças de proporções elevadas podem ser produzidas por um número muito pequeno de observações. Por exemplo, uma proporção de 100\% pode corresponder tanto a uma ocorrência em uma única observação quanto a um padrão repetido em um número elevado de implementações.

Como diagnóstico exploratório adicional, considera-se o suporte mínimo:

\[
S=\min(n_X,n_Y).
\]

Essa quantidade representa a menor frequência marginal entre a alucinação e a regra analisadas. O suporte mínimo não constitui uma medida de associação nem um critério estatístico de significância. Ele é utilizado apenas para auxiliar a interpretação da magnitude dos padrões observados.

A relação entre $\Delta$ e as frequências $n_X$, $n_Y$ e $S$ é examinada graficamente para verificar se diferenças extremas estão concentradas principalmente em pares com baixo suporte observacional.

Não é estabelecido, a priori, um ponto de corte para classificar um par como suficientemente suportado. Em vez disso, magnitude, direção e suporte observacional são considerados conjuntamente na interpretação dos padrões.

\subsubsection{Identificação dos padrões exploratórios}

A caracterização de cada par alucinação--regra considera conjuntamente:

\[
a,\quad
n_X,\quad
n_Y,\quad
P(Y\mid X),\quad
P(Y\mid \neg X),\quad
\Delta
\]

e, quando informativo, o \textit{odds ratio}.

A análise busca identificar:

\begin{itemize}
    \item padrões positivos, nos quais a ocorrência de $Y$ é proporcionalmente maior entre implementações com $X$;
    \item padrões negativos, nos quais a ocorrência de $Y$ é proporcionalmente menor entre implementações com $X$;
    \item casos de ausência de coocorrência em que $X$ e $Y$ apresentam frequências marginais suficientes para que esse padrão seja descritivamente relevante;
    \item valores extremos associados a frequências muito baixas, que devem ser interpretados com cautela.
\end{itemize}

Dessa forma, nenhum par é caracterizado como relevante exclusivamente com base em um valor elevado de $\Delta$ ou do OR. A interpretação considera simultaneamente a direção, a magnitude e o suporte observacional do padrão.

\subsection{Análise no nível da tarefa}

Para a análise no nível da tarefa, os indicadores binários das cinco implementações correspondentes a uma mesma tarefa são agregados utilizando o critério \textit{any}.

Para cada alucinação $X$:

\[
X_{\text{task}}=
\begin{cases}
1, & \text{se pelo menos uma implementação da tarefa apresenta }X;\\
0, & \text{caso contrário.}
\end{cases}
\]

Analogamente, para cada regra $Y$:

\[
Y_{\text{task}}=
\begin{cases}
1, & \text{se pelo menos uma implementação da tarefa apresenta }Y;\\
0, & \text{caso contrário.}
\end{cases}
\]

A análise exploratória no nível da tarefa seguirá a mesma lógica geral utilizada no nível da implementação, com construção de tabelas de contingência e avaliação das proporções, magnitude e suporte dos padrões observados.

A definição final das medidas e visualizações utilizadas nesta etapa será estabelecida após a conclusão da análise exploratória no nível da implementação.